% =====================================================
% CHAPITRE 2 : COMPRÉHENSION MÉTIER (CRISP-DM Phase 1)
% =====================================================
\chapter{Compréhension métier}

% ---------------------------------------------------
\section*{Introduction}
\addcontentsline{toc}{section}{Introduction}
% ---------------------------------------------------

Avant d'écrire la moindre ligne de code, il faut comprendre le problème que l'on cherche
à résoudre --- et surtout, le métier derrière les données. C'est l'essence de la première
phase CRISP-DM. Ce chapitre pose donc le cadre : contexte de la fraude financière telle
qu'on la rencontre en mission d'audit, formalisation de la problématique, objectifs
opérationnels et critères de succès mesurables qui guideront toutes les décisions techniques
ultérieures. Cette étape est non négociable dans un contexte d'audit IT : une définition
floue des objectifs aurait conduit à des choix algorithmiques déconnectés des besoins réels
des auditeurs.

% ---------------------------------------------------
\section{Contexte de la détection de fraude financière}
% ---------------------------------------------------

\subsection{Enjeux économiques et opérationnels}

La fraude financière représente une menace persistante et coûteuse pour les institutions.
Si les chiffres globaux donnent le vertige --- le rapport Nilson (2022) évalue à plusieurs
dizaines de milliards de dollars les pertes annuelles liées à la fraude par carte bancaire
---, c'est à l'échelle d'une mission d'audit que l'impact se fait le plus ressentir :
transactions non reconnues, coûts de remboursement, investigations manuelles chronophages
et responsabilité engagée des auditeurs. À cela s'ajoutent des contraintes réglementaires
croissantes, comme la directive européenne PSD2 qui impose des mécanismes de surveillance
renforcés des transactions.

Pour une équipe comme celle du Risk Assurance Services de PwC Tunisie, l'enjeu est double :
détecter les anomalies avant qu'elles ne se propagent, et le faire de manière traçable et
défendable dans un rapport d'audit.

\subsection{Limites des systèmes basés sur des règles}

Historiquement, la détection de fraude repose sur des règles métier définies manuellement
par des experts : seuils de montant, listes noires de comptes suspects, patterns connus de
détournement. Ces systèmes ont le mérite d'être explicables et rapides à déployer. Mais
lors de l'analyse du corpus de ce stage, le constat a été frappant : le système naïf
\texttt{isFlaggedFraud} présent dans les données n'identifie qu'un seul cas de fraude sur
258, soit un rappel de \textbf{0,39\,\%}. Trois limites fondamentales expliquent cet
échec :

\begin{itemize}
  \item \textbf{Rigidité} : les règles codifient les comportements connus, mais les
    fraudeurs adaptent continuellement leurs méthodes pour les contourner.
  \item \textbf{Maintenance coûteuse} : chaque nouvelle forme de fraude nécessite une
    intervention manuelle d'un expert pour mettre à jour les règles.
  \item \textbf{Faux positifs élevés} : les seuils statiques génèrent de nombreuses
    alertes injustifiées, surchargeant les équipes d'investigation et dégradant la
    confiance dans le système.
\end{itemize}

L'apprentissage automatique offre une alternative complémentaire : plutôt que d'encoder les
patterns de fraude à la main, il les extrait automatiquement des données historiques, y
compris des interactions non linéaires entre variables que l'œil humain ne perçoit pas.

% ---------------------------------------------------
\section{Présentation et formalisation du projet}
% ---------------------------------------------------

\subsection{Contexte du projet de fin d'études}

Ce projet a été réalisé dans le cadre d'un stage de fin d'études axé sur la mise en œuvre
d'un système de détection d'anomalies financières exploitant des techniques avancées
d'apprentissage automatique. Le travail couvre l'intégralité du cycle de vie d'un projet
data science, de la définition du problème au déploiement d'un pipeline reproductible.

\subsection{Problématique}

\begin{definitionbox}
\textbf{Problématique :} Comment concevoir un pipeline reproductible, fondé sur la
méthodologie CRISP-DM, capable de détecter efficacement les transactions financières
frauduleuses dans un contexte de fort déséquilibre des classes, tout en produisant des
explications compréhensibles pour les analystes métier ?
\end{definitionbox}

Cette problématique est structurée autour de deux défis techniques majeurs : (1) la
\textbf{gestion du déséquilibre extrême} des classes --- moins de 0,13\,\% des transactions
sont frauduleuses, un ratio si faible qu'un modèle naïf prédisant systématiquement
« légitime » afficherait 99,87\,\% de précision tout en ne détectant aucune fraude ---
et (2) l'\textbf{interprétabilité} des décisions, rendue possible grâce aux méthodes SHAP,
LIME et à la génération d'explications textuelles par le modèle de langage Ollama.

\subsection{Objectifs du projet}

Les objectifs opérationnels définis en début de projet sont les suivants :

\begin{enumerate}
  \item Construire une \textbf{chaîne de traitement complète et reproductible} couvrant
    toutes les phases CRISP-DM.
  \item Réaliser une \textbf{analyse exploratoire approfondie} des données transactionnelles.
  \item Concevoir et appliquer un \textbf{pipeline de préparation des données} robuste,
    préservant l'intégrité temporelle et évitant toute fuite d'information
    (\textit{data leakage}).
  \item \textbf{Entraîner et comparer} au moins trois approches de détection : baseline
    linéaire, modèle ensembliste supervisé et autoencoder non supervisé.
  \item Évaluer les modèles avec des \textbf{métriques adaptées au déséquilibre} : rappel,
    précision, F1-score, courbe PR et PR-AUC.
  \item Produire des \textbf{explications interprétables} des décisions (SHAP, LIME, Ollama).
  \item Livrer un \textbf{code modulaire, documenté et testé} facilitant la maintenance
    et l'extension.
\end{enumerate}

\subsection{Critères de succès}

\begin{table}[H]
\centering
\caption{Critères de succès du projet}
\label{tab:criteres_succes}
\begin{tabular}{lll}
\toprule
\textbf{Critère} & \textbf{Valeur cible} & \textbf{Justification} \\
\midrule
Recall (test)    & $> 0{,}70$ & Priorité à la détection des fraudes \\
F1-score (test)  & $> 0{,}70$ & Équilibre précision/rappel \\
PR-AUC (test)    & $> 0{,}75$ & Synthèse du tradeoff \\
Reproductibilité & 100\,\%   & Pipeline exécutable via \texttt{run\_all.py} \\
Couverture tests & $> 80\,\%$ & Robustesse du code \\
\bottomrule
\end{tabular}
\end{table}

% ---------------------------------------------------
\section{Cadre méthodologique : CRISP-DM}
% ---------------------------------------------------

\subsection{Présentation de la méthodologie}

\begin{definitionbox}
\textbf{CRISP-DM} (\textit{Cross-Industry Standard Process for Data Mining}) est un cadre
méthodologique ouvert proposé en 1996 par un consortium réunissant SPSS, Daimler-Benz,
NCR et OHRA. Son objectif : structurer de façon reproductible les projets de fouille de
données, quelle que soit l'industrie ou les outils utilisés. Il organise le travail en
\textbf{six phases interdépendantes} formant un cycle itératif, du cadrage métier jusqu'à
la mise en production.
\end{definitionbox}

Les six phases de CRISP-DM sont :

\begin{enumerate}
  \item \textbf{Business Understanding} --- Compréhension métier et définition des objectifs.
  \item \textbf{Data Understanding} --- Collecte, exploration et évaluation de la qualité
    des données.
  \item \textbf{Data Preparation} --- Sélection, nettoyage, construction et formatage des
    données.
  \item \textbf{Modeling} --- Sélection, construction et ajustement des modèles.
  \item \textbf{Evaluation} --- Évaluation des résultats au regard des objectifs métier.
  \item \textbf{Deployment} --- Mise en production et documentation.
\end{enumerate}

\subsection{Justification du choix de CRISP-DM}

Plusieurs méthodologies structurées existent pour conduire des projets de data science.
Le tableau ci-dessous les compare pour justifier le choix retenu.

\begin{table}[H]
\centering
\caption{Comparaison des principales méthodologies de data science}
\label{tab:methodo_compare}
\begin{tabular}{p{2.3cm}p{3.4cm}p{3.4cm}p{3.2cm}}
\toprule
\textbf{Méthodologie} & \textbf{Points forts} & \textbf{Limites}
  & \textbf{Adéquation} \\
\midrule
\textbf{CRISP-DM} &
  Agnostique outils ; cycle itératif ; couverture complète (métier $\to$ déploiement) &
  Moins adapté aux projets agiles purs &
  \textbf{Excellent} \\[0.3cm]
\textbf{SEMMA} (SAS) &
  Simple et concis &
  Lié à SAS ; absence de phase déploiement &
  \textbf{Insuffisant} \\[0.3cm]
\textbf{KDD} &
  Fondement académique solide &
  Peu flexible ; orienté recherche &
  \textbf{Partiel} \\[0.3cm]
\textbf{Agile DS} &
  Itérations courtes, DevOps &
  Suppose équipe multidisciplinaire et CI/CD &
  \textbf{Inadapté} \\
\bottomrule
\end{tabular}
\end{table}

Le choix de CRISP-DM pour ce projet n'était pas arbitraire. Plusieurs raisons concrètes
l'ont motivé :

\begin{itemize}
  \item \textbf{Orientation métier explicite} : avant de toucher aux données, CRISP-DM
    impose de formaliser les objectifs et les critères de succès --- étape non négociable
    en contexte d'audit IT.
  \item \textbf{Couverture complète du cycle} : contrairement à SEMMA ou KDD, CRISP-DM
    couvre aussi bien la compréhension des données que le déploiement du pipeline.
  \item \textbf{Caractère itératif} : la préparation des données et la modélisation ne
    sont pas séquentielles en pratique. Les allers-retours entre phases sont formalisés
    et légitimés par CRISP-DM.
  \item \textbf{Indépendance des outils} : CRISP-DM étant agnostique, il s'adapte sans
    difficulté à une stack Python / scikit-learn / PyTorch / Ollama.
  \item \textbf{Lisibilité pour les équipes PwC} : CRISP-DM étant largement adopté dans
    les cabinets de conseil, les encadrants entreprise ont pu suivre la progression du
    projet sans avoir à maîtriser les détails techniques.
\end{itemize}

\subsection{Application au projet}

L'ensemble du rapport est structuré selon ces six phases. Le caractère itératif de
CRISP-DM se manifeste notamment dans les allers-retours entre l'analyse exploratoire
(phase 2) et la préparation des données (phase 3), ainsi qu'entre la modélisation
(phase 4) et l'évaluation (phase 5) lors du réglage des seuils de détection.

% ---------------------------------------------------
\section*{Conclusion}
\addcontentsline{toc}{section}{Conclusion}
% ---------------------------------------------------

Ce chapitre a établi le contexte métier du projet, la problématique à résoudre et les
objectifs à atteindre selon le cadre CRISP-DM. Deux défis principaux guideront toutes
les décisions techniques : la gestion du déséquilibre extrême des classes (ratio 1:774)
et la nécessité de produire des décisions interprétables pour les auditeurs. Le choix
de CRISP-DM comme méthodologie directrice a été justifié par sa couverture complète du
cycle de vie, sa nature itérative et son indépendance vis-à-vis des outils. Le chapitre
suivant présente les données sur lesquelles repose l'intégralité du travail.
